\section{Threats to Validity}
\label{sec:threats}

\textbf{Conclusion validity.} The H2 rejection margin ($+0.0090$
vs $+0.01$ pre-registered) is within the 95\% bootstrap CI for the
cell-mean comparison. We report the binary rejection per the
protocol and the substantive direction in the abstract.

\textbf{Internal validity.}

\textit{(i) CUSUM $(k, h)$ is pre-registered.} The $(0.04, 0.10)$
values were locked before any cusum P@50 was inspected. They were
chosen by inspection of Paper~11's calibration-seed cp\_delta
distribution: $k = 0.04$ is half the typical K=50 |Delta|; $h =
0.10$ is the largest single-$\tau$ Paper~11 swept. A different
$(k, h)$ pair would shift the K=200 gain and K=50 cost in
correlated ways; the present paper measures one point on this
two-parameter family.

\textit{(ii) Single $\lambda$ and capacity grid.} As in
Papers~7--12. Different $\lambda$ would shift the per-window
$|\Delta_w|$ distribution and re-locate the CUSUM operating point.

\textit{(iii) Selection-policy coupling.} Inherited from Papers~7--12.

\textbf{Construct validity.} CUSUM's accumulator state persists
across windows within a cell; the strategy's behaviour is therefore
\textit{path-dependent} in a way that single-window detectors are
not. The K=200 gain depends on the specific pattern of
$|\Delta_w|$ values, including the W4 noise excursion that the
synthetic generator produces. A real fleet with a different
calibration-target time series would produce different fire
decisions; the qualitative claim (CUSUM filters single-window noise)
is more robust than the specific $+0.9$~pp gain magnitude.

\textbf{External validity.} EEHDA synthetic priors inherited from
Papers~4--12. The K=200 W4 noise excursion is a property of the
fleet evolution model; real fleet calibration-target excursions
may have different temporal structure.
